% !TEX root = ../main.tex
\chapter{The Syntax and Semantics of Generalized Algebraic Theories}

In this chapter we will introduce Generalized Algebraic Theories, a framework for 
working with dependent types. They allow us to construct the concept of contextual 
category, a key base structure for interpreting Calculus of Constructions.


\section{Generalized Algebraic Theories}

Generalized Algebraic Theory (or GAT), originally presented by J. Cartmell 
in his thesis \cite{cartmell1978} \cite{cartmell1986},
is an abstraction of Martin-L\"of type theory
and is a generalization of the usual notion of a many-sorted 
algebraic theory in the following manner:
while in Martin L{\"o}ff's many-sorted algebraic theories sorts
are constant types, interpreted as sets, in generalized
algebraic theories sorts might be constant or variable,
over a specified range of types, and variable types
are interpreted as families of sets. So a generalized
algebraic theory consists of:
\begin{enumerate}
    \item a set of sorts, that can be constant or variable;
    \item a set of operators, with the sort of the arguments and value specified;
    \item a set of axioms, where each axiom is an identity between similar expressions. 
\end{enumerate}

Let's present category theory for example:
The set of sorts contains two elements $Ob$ and $Hom$.
$Ob$ is a constant type and $Hom$ is a variable type,
for a pair of terms of type $Ob$, name them $x, y$, we
have a type $Hom(x, y)$.
Here we see a key point in general algebraic theories,
types vary over terms, no over types.
The set of operators also has two elements $id$ and $\circ$.
The $id$ operator has one argument $x$ of type $Ob$,
and returns a value of type $Hom(x, x)$. This shows
how the type of value of an operator might vary depending
on the arguments. 
The $\circ$ operator (or composition) expects 5 arguments, 
three of type $Ob$, name them $x, y, z$, one of type $Hom(x, y)$,
say $f$, and one of type $Hom(y, z)$, named $g$, and it returns
a value of type $Hom(x, z)$.
Here we see that the type of an argument might also vary depending
on a previous argument.
Informally, instead of writing $\circ(x, y, z, f, g)$ we can
write $\circ(f, g)$, as the first three arguments might be inferred from the last two.
This notation is discussed in the original paper \cite{cartmell1986}. 

Now we will introduce the notation we will be using when talking 
about terms and types. {\color{orange} This notation will be formally defined
in a later section.} So consider you have a set of variables $V$,
we will associate terms to variables as required.
When asserting that a term $t$ is of type $A$ we will write
$t : \Delta$, but if $t$ has variables $x_1, \dots, x_n$ occurring in it this assertion
does not make sense unless we assign some types to $x_1, \dots, x_m$.
To do this we will write:
\begin{equation*}
    \dfrac{
        x_1 : \Delta_1, \dots, x_n: \Delta_n
    }{
        t : \Delta
    }
\end{equation*}
that reads as: if $x_1$ is a variable of type $\Delta_1$, \dots and if $x_n$ is a variable of type $\Delta_1$,
then $t$ is a term of type $\Delta$. 
Similarly, if we want to assert that if $x_1$ is a variable of type $\Delta_1$,
\dots and if $x_n$ is a variable of type $\Delta_n$,
then $\Delta$ is a type, we can write:
\begin{equation*}
    \dfrac{
        x_1 : \Delta_1, \dots, x_n : \Delta_n 
    }{
        t : \Delta
    }
\end{equation*}
this assertions are called rules or sequents.
Rules are used to express which expressions are well-formed.
We will work with this rules as units instead of the base expressions.
This is because terms and types do not make sense by themselves
but they need to be thought of in a context. 

Axioms are also written as rules. The conclusion of those rules
are equalities between types or between terms of the same type.
We write those rules as:
\begin{minipage}{0.5\textwidth}
  \begin{equation*}
    \dfrac{
        x_1 : \Delta_1, \dots, x_n : \Delta_n
    }{
        \Delta = \Delta'
    }
  \end{equation*}
\end{minipage}%
\begin{minipage}{0.5\textwidth}
  \begin{equation*}
    \dfrac{
        x_1 : \Delta_1, \dots, x_n : \Delta_n
    }{
        t = t' : \Delta
    }
  \end{equation*}
\end{minipage}
and read as: if $x_1$ is a variable of type $\Delta_1$, \dots and $x_n$ is a variable of type $\Delta_n$,
then $\Delta = \Delta'$, or $t = t'$ as terms of type $\Delta$, respectively.

% TODO: fix alignment
To continue with the previous example, we can write the axioms
of category theory as:
\begin{alignat*}{2}
    &x, y : Ob, f : Hom(x, y),  \quad  \circ(id(x), f) = f \\
    &x, y : Ob, f : Hom(x, y),  \quad  \circ(f, id(y)) = f \\
    &w, x, y, z : Ob, f : Hom(w, x), g : Hom(x, y), h : Hom(y, z), \\ &\qquad\circ(\circ(f, g), h) = \circ(f, \circ(g, h))
\end{alignat*}

When presenting a theory's language i.e. the sets of sorts and operators,
we must present each symbol with a rule specifying the role
of this symbol, either as a sort symbol or as specifically typed operator symbol. 
This rules are called introductory rules.
In the case of a sort symbol $A$ the rule should be like 
\begin{equation*}
    \dfrac{
        x_1 : \Delta_1, \dots, x_n : \Delta_n
    }{
        A(x_1, \dots, x_n) \text{ is a type}
    }
\end{equation*}
clearly specifying the types $A$ is dependent on.
In the case of an operator symbol, a rule of the form
\begin{equation*}
    \dfrac{
        x_1 : \Delta_1, \dots, x_n : \Delta_n
    }{
        f(x_1, \dots, x_n) : \Delta
    }
\end{equation*}
suffices to specify the types of the arguments and the value of $f$.

Finally, to present a theory we will give a set os symbols,
each associated to its introductory rule, and a set of axioms.
And, of course, everything must be well-formed, {\color{orange} something we 
will address in the following section.}
So now we can give a formal presentation of the theory of categories:

\begin{tabular}{ll}
\textit{Symbol} & \textit{Introductory Rule} \\[6pt]
Ob  & Ob is a type. \\[4pt]
Hom & $x, y \in \mathrm{Ob} : \mathrm{Hom}(x, y)$ is a type. \\[4pt]
$\circ$ & $x, y, z \in \mathrm{Ob},\, f \in \mathrm{Hom}(x, y),\, g \in \mathrm{Hom}(y, z) : \circ(f, g) \in \mathrm{Hom}(x, z).$ \\[4pt]
id  & $x \in \mathrm{Ob} : \mathrm{id}(x) \in \mathrm{Hom}(x, x).$ \\
\end{tabular}

\begin{tabular}{ll}
\textit{Axioms} \\[6pt]
$\circ(\mathrm{id}(x), f) = f,$ & whenever $x, y \in \mathrm{Ob}$ and $f \in \mathrm{Hom}(x, y).$ \\[4pt]
$\circ(f, \mathrm{id}(y)) = f,$ & whenever $x, y \in \mathrm{Ob}$ and $f \in \mathrm{Hom}(x, y).$ \\[4pt]
$\circ(\circ(f, g), h) = \circ(f, \circ(g, h)),$ 
  & whenever $w, x, y, z \in \mathrm{Ob},\, f \in \mathrm{Hom}(w, x),$ \\
  & \hspace{4.5em} $g \in \mathrm{Hom}(x, y)$ and $h \in \mathrm{Hom}(y, z).$ \\
\end{tabular}

\section{Formal Definition}

In this section we will formally define all of the concepts 
we just introduced.
To give a notion of well-formed introductory rule, we need the concept of drivability,
which also depends on the introductory rules. 
To tackle this problem we will leave aside the question of
well-formedness until we have the complete set of derived rules
of a theory.
We will first define a pretheory and then define a theory
as a well-formed pretheory.

We will assume to have countable set of variables $V$.
Lets define the notion of rule on an alphabet $W$,
the elements of $W$ are called symbols.

\begin{definition}
    The set of \textbf{expressions} is defined inductively
    such that expressions are finite sequences of elements
    of $W \cup V \cup \{(\} \cup \{)\} \cup \{,\}$
    defined by the clauses:
    \begin{enumerate}
        \item if $x \in V$, then $x$ is an expression;
        \item if $L \in W$, then $L$ is an expression;
        \item if $L \in W$ and $e_1, \dots, e_n$ are expressions,
              then $L(e_1, \dots, e_n)$ is an expression.
    \end{enumerate}
\end{definition}

Now we can define the building blocks of a rule:
\begin{definition}
    A \textbf{premise} is defined to be a sequence of $V \times$ the set of expressions.
    The empty sequence is allowed and we call it an empty premise.
    If $((x_1, \Delta_1), \dots, (x_n, \Delta_n))$ is a premise,
    we write it as $x_1 : \Delta_1, \dots, x_n : \Delta_n$.
\end{definition}
\begin{definition}
    We call a \textbf{conclusion} to one of the following:
    \begin{enumerate}
        \item $T$-conclusion: defined by a single expression $\Delta$
              and written as '$\Delta$ is a type'.
        \item $\in$-conclusion: defined by a pair of expressions $(t, \Delta)$
              and written as '$t \in \Delta$'.
        \item $T=$conclusion: defined by a pair of expressions $(\Delta_1, \Delta_2)$
              and written as '$\Delta_1 = \Delta_2$'.
        \item $\in=$conclusion: defined by a triple of expressions $(t_1, t_2, \Delta)$
              and written as '$t_1 = t_2 : \Delta$'
    \end{enumerate}
\end{definition}
Finally a \textbf{rule} is determined by a premise $P$ and a
conclusion $C$, and is written as $\dfrac{P}{C}$.
We name the rule as $T$-rule, $\in$-rule, $T=$rule or $\in=$rule
according to the form of the conclusion.

Now we can define pretheory and later theory.

\begin{definition}
    A \textbf{pretheory} consists of
    \begin{enumerate}
        \item a set of sort symbols $S$;
        \item a set of operator symbols $Z$;
        \item for each sort symbol $A$, an associated rule 
              of the alphabet $S \cup Z$:
              \begin{equation*}
                \dfrac{
                    x_1 : \Delta_1, \dots, x_n : \Delta_n
                }{
                    A(x_1, \dots, A_n) \text{ is a type}
                }
              \end{equation*}
              for some $n \geq 0$. It is called the introductory
              rule of $A$;
        \item for each operator symbol $F$, an associated rule
              of the alphabet $S \cup Z$:
              \begin{equation*}
                \dfrac{
                    x_1 : \Delta_1, \dots, x_n : \Delta_n
                }{
                    F(x_1, \dots, x_n) : \Delta
                }
              \end{equation*}
              for some $n \geq 0$. It is called the introductory
              rule of $F$;
        \item a set of axioms. Each axiom is either a $T=$rule or
              a $\in=$rule of the alphabet $S \cup Z$.

    \end{enumerate}
\end{definition}

Now we will tackle the problem of well-formedness of rules.

\begin{definition}
    If $U$ is a pretheory, then
    \begin{enumerate}
        \item a \textbf{context} is a premise $x_1 : \Delta_1, \dots, x_n : \Delta_n$
            such that the rule
            \begin{equation*}
                \dfrac{
                    x_1 : \Delta_1, \dots, x_{n-1}: \Delta_{n-1}
                }{
                    \Delta_n \text{ is a type}
                }
            \end{equation*}
            is a derived rule such that $x_n$ is distinct from all of $x_1, \dots, x_n$.
        \item \begin{enumerate}
            \item the rule $\dfrac{x_1 : \Delta_1, \dots, x_n : \Delta_n}{\Delta \text{ is a type}}$
                  is \textbf{well-formed} if the premise $x_1 : \Delta_1, \dots, x_n : \Delta_n$ is a context.
            \item the rule $\dfrac{x_1 : \Delta_1, \dots, x_n : \Delta_n}{t : \Delta}$
                  is \textbf{well-formed} if $\dfrac{x_1 : \Delta_1, \dots, x_n : \Delta_n}{\Delta \text{ is a type}}$
                  is a derived rule of $U$.
            \item the rule $\dfrac{x_1 : \Delta_1, \dots, x_n : \Delta_n}{\Delta = \Delta'}$
                  is \textbf{well-formed} if $\dfrac{x_1 : \Delta_1, \dots, x_n : \Delta_n}{\Delta \text{ is a type}}$
                  and \\ $\dfrac{x_1 : \Delta_1, \dots, x_n : \Delta_n}{\Delta' \text{ is a type}}$
                  are derived rules of $U$.
            \item the rule $\dfrac{x_1 : \Delta_1, \dots, x_n : \Delta_n}{t = t' : \Delta}$
                  is \textbf{well-formed} if $\dfrac{x_1 : \Delta_1, \dots, x_n : \Delta_n}{t : \Delta}$
                  and \\ $\dfrac{x_1 : \Delta_1, \dots, x_n : \Delta_n}{t' : \Delta}$
                  are derived rules of $U$.
        \end{enumerate}
    \end{enumerate}
\end{definition}

\begin{definition}
    The set of \textbf{derived rules} of a pretheory $U$ is the
    set of rules derivable by the following principles of derivation:
    
    % TODO: add rules of derivation
\end{definition}

\begin{definition}
    A pretheory is \text{well-formed} if all of its introductory
    rules are well-formed.
    A \textbf{theory} is a well-formed pretheory.
\end{definition}

Lastly, we will introduce substitution.
Given $\Delta, t_1, \dots, t_n$ expressions and \\ 
$x_1, \dots, x_n$ variables,
then we define the expression $\Delta[x_1 | t_1, \dots, x_n | t_n]$
the result of simultaneous replacement of $x_1, \dots, x_n$ in $\Delta$
by $t_1, \dots, t_n$. We call this the Substitution operation.
Then, as proven in the original paper \cite{cartmell1986}, 
the set of derived rules of a theory is closed under the 
operation of substitution of correctly typed terms 
for variables.
We could add substitution as one of the derivation rules
but would make applying induction very messy.

\section{Models and Homomorphisms} % TODO: probably useless

With all the syntax defined, we will consider the semantics of a theory.
To do this, we will introduce the notions of model and model homomorphism.

First, we will set some notation and basic concepts.

% TODO: Delete this as it defined elsewhere ----
\begin{definition}
    A \text{tree} is a pair $(N, father)$ where 
    $N$ is a not empty set of nodes and 
    $father: N \to N$ such that:
    \begin{enumerate}
        \item for all $A \in N$ there exists $n \in \omega$
              such that $father^n(A) = father^{n+1}(A)$;
        \item \label{item:tree-ii.aux}
              for all $A, B \in N$, if $A = father(A)$ and
              $B = father(B)$, then $A = B$.
    \end{enumerate}
\end{definition}
We will write $A \unlhd B$ if $A = father(B)$
and $A \lhd B$ if $A \unlhd B$ and $A \neq B$.
By condition \ref{item:tree-ii.aux}, there is a unique
least element of $N$, call it $1$.
Finally we will define, for any $A \in N$,
$level(A)$ to be the least $n \in \omega$ with 
$father^n(A) = father^{n+1}(A)$.

The set of contexts of a theory is structured as a tree.
Whit the empty context as it's least element and 
\(
    father (\langle x_1 : \Delta_1, \dots, x_n : \Delta_n\rangle)
    = 
    \langle x_1 : \Delta_1, \dots, \Delta_{n-1} \rangle
\).

% TODO: ----

We will also want to consider the tree of sets, families of sets, families of families of sets, \dots
Consider the following notation for families:
given a set $A$, if for every element $a \in A$ $B(A)$ is a set,
then the corresponding $A$-indexed family of sets is denoted as
$\lambda a \in A . B(a)$ or as $\lambda a . B(a)$;
If we have a set $C(a, b)$ for every $a \in A$ and $b \in B(a)$,
then $\lambda a \in A. \lambda b \in B(a). C(a, b)$
is an $A$-indexed family of families of sets; and so on.
The collection of sets, families of sets, families of families of sets, and so on,
is called the \textbf{tree of families}.
Given $1 \vartriangleleft A_1 \vartriangleleft A_2 \vartriangleleft \dots \vartriangleleft A_n \vartriangleleft A$
in the tree of families. Then, for all $a_1 \in A_1$, \dots,
for all $a_n \in A_n(a_1, \dots, a_{n-1})$, $A(a_1, \dots, a_n)$ is a set.
Finally we should define the concept of \textbf{operator}.
If we have $1 \vartriangleleft A_1 \vartriangleleft \dots \vartriangleleft A_n \vartriangleleft A$
in the tree of families. If for all $a_1 \in A_1$, \dots, 
for all $a_n \in A_n(a_1, \dots, a_{n-1})$,
$f(a_1, \dots, a_n) \in A(a_1, \dots, a_n)$. 
We say that $\lambda a_1 \in A_1. \dots \lambda a_n \in A_n(a_1, \dots, a_{n-1}). f(a_1, \dots, a_n)$
is an operator at $A$.

Now we have the necessary structure to define a model of a theory. It is necessary to state the fact that 
models do not interpret expressions, they interpret the derived rules of the theory. The interpretation
of a rule $R$ within a model $M$ is written as $R^M$.

\begin{definition}
    Let $U$ be a theory, a \textbf{model} $M$ is a partial mapping from the set of derived rules of $U$ to $Fam$,
    such that:
    \begin{enumerate}
        \item if $R$ is a derived $T$-rule \(
            \dfrac{
                x_1 \colon \Delta_1, \dots, x_n \colon \Delta_n
            }{
                \Delta \text{ is a type}
            }
        \), if we name $R_i$ the rules \(
            \dfrac{
                x_1 \colon \Delta_1, \dots, x_{i-1} \colon \Delta_{i-1}
            }{
                \Delta_i \text{ is a type}
            }
        \) for $1 \leq i \leq n$. 
        Then $1 \lhd R^M_1 \lhd R^M_2 \lhd \dots \lhd R^M_n \lhd R^M$ in the tree of families.
        \item if $R_t$ is a derived $\in$-rule \(
            \dfrac{
                x_1 \colon \Delta_1, \dots, x_n \colon \Delta_n
            }{
                t \colon \Delta
            }
        \) the interpretation $R^M_t$ is an operator at $R^M$.
    \end{enumerate}
\end{definition}

% TODO: add definition of homomorphism ?

\section{Contexts and Realizations} % --- Notation Changes ---

In order to construct the contextual category of a theory $U$,
we will first define the category of contexts and realizations of $U$.
This category, noted $R(U)$ has as objects the contexts of $U$.
Given two contexts $\Gamma \equiv x_1 \colon A_1, \dots, x_n \colon A_n$ and 
$\Delta \equiv y_1 \colon B_1, \dots, y_m \colon B_m$ of $U$, a morphism from $\Gamma$ to $\Delta$,
or a \textit{realization} of $\Delta$ from $\Gamma$, is a $m$-tuple $\langle t_1, \dots, t_n \rangle$ 
such that for each $j$, $1 \leq j \leq m$, the rule 
\[
    \vdash \Gamma \Rightarrow t_j \colon B_j[y_1 \mid\ t_1, \dots, y_{j-1} \mid\ t_{j-1}] 
\]
The identity of an object $x_1 \colon \Delta_1, \dots, x_n \colon \Delta_n$ is the realization 
$\langle x_1, \dots, x_n \rangle$, and composition is defined as 
\[
    \langle t_1, \dots, t_m \rangle \circ \langle s_1, \dots, s_w \rangle =
    \langle s_1[y_1 \mid\ t_1, \dots, y_m \mid\ t_m], \dots, s_w[y_1 \mid\ t_1, \dots, y_m \mid\ t_m]
\]
whenever $\Gamma = x_1 \colon A_1, \dots, x_n \colon A_n$, $\Delta = y_1 \colon B_1, \dots, y_m \colon B_m$
and $\Omega = z_1 \colon C_1, \dots,  z_w \colon C_w$ are contexts, and 
$\langle t_1, \dots, t_m \rangle : \Gamma \to \Delta$ and $\langle s_1, \dots, s_w \rangle : \Delta \to \Omega$
are realizations.

This category has two important properties we should emphasize in:

Given the following definition of tree:
\begin{definition}
    A \text{tree} is a pair $(N, father)$ where 
    $N$ is a not empty set of nodes and 
    $father: N \to N$ such that:
    \begin{enumerate}
        \item for all $A \in N$ there exists $n \in \omega$
              such that $father^n(A) = father^{n+1}(A)$;
        \item \label{item:tree-ii}
              for all $A, B \in N$, if $A = father(A)$ and
              $B = father(B)$, then $A = B$.
    \end{enumerate}
\end{definition}
We will write $A \unlhd B$ if $A = father(B)$
and $A \lhd B$ if $A \unlhd B$ and $A \neq B$.
By condition \ref{item:tree-ii}, there is a unique
least element of $N$, call it $1$.
Finally we will define, for any $A \in N$,
$level(A)$ to be the least $n \in \omega$ with 
$father^n(A) = father^{n+1}(A)$.
We can see that the set of objects of $R(U)$, i.e. the set of contexts of $U$ is structured as a tree
Whit the empty context as it's least element and 
\(
    father (\langle x_1 : \Delta_1, \dots, x_n : \Delta_n\rangle)
    = 
    \langle x_1 : \Delta_1, \dots, \Delta_{n-1} \rangle
\).
Moreover, given a context $x_1 \colon A_1, \dots, x_n \colon A_n, x_{n+1} \colon A_{n+1}, \dots x_{n+m}$
of $U$, $\langle x_1, \dots, x_n \rangle$ is a realization of $x_1 \colon A_1, \dots, x_n \colon A_n$
with regard to $x_1 \colon A_1, \dots, x_n \colon A_n, x_{n+1} \colon A_{n+1}, \dots x_{n+m}$, and is noted
as $p(\langle x_1 \colon A_1, \dots, x_{n+m} \colon A_{n+m} \rangle, \langle x_1 \colon A_1, \dots, x_n \colon A_n \rangle)$
Hence, given two objects $\Gamma, \Delta \in \Ob R(U)$, such that $\Gamma \leq \Delta$, $p(\Delta, \Gamma)$ is defined and 
$p(\Delta, \Gamma) : \Delta \to \Gamma$ in $R(U)$.

And if $\Gamma \equiv x_1 \colon A_1, \dots, x_n \colon A_n$, $\Gamma' \equiv y_1 \colon B_1, \dots, y_m \colon B_m$
and $\Delta \equiv y_1 \colon B_1, \dots, y_{m+w} \colon B_{m+w}$ are contexts, and if 
$f = \langle t_1, \dots, t_m \rangle : \Gamma \to \Gamma'$ in $R(U)$, then
\[
\begin{tikzcd}
    f^*\Delta \arrow[r, "{q(f, \Delta)}"] \arrow[d, "{p(f^*\Delta, \Gamma)}"'] & \Delta \arrow[d, "{p(\Delta, \Gamma')}"] \\
    \Gamma \arrow[r, "f"] & \Gamma'
\end{tikzcd}
\]
is a pullback diagram in $R(U)$, where % TODO: fix alignment
\begin{multline*}
    f^*\Delta 
    \equiv 
        x_1 \colon A_1, \dots, x_n \colon A_n, \\
        y_{m+1} \colon B[y_1 \mid\ t_1, \dots, y_m \mid\ t_m], 
        \dots,
        y_{m+w} \colon B[y_1 \mid\ t_1, \dots, y_m \mid\ t_m] 
\end{multline*}
\[
    q(f, \Delta) = \langle t_1, \dots, t_m, y_{m+1}, \dots, y_{m+w} \rangle
\]

This motivates the following definition

\section{Contextual Categories}

\begin{definition}
    a \textbf{contextual category} consists of
    \begin{enumerate}
        \item a category $\mathcal{C}$ with a terminal object $1$;
        \item a function $\father \colon \Ob \mathcal{C} \to \Ob \mathcal{C}$ 
              such that $(\Ob \mathcal{C}, \father)$ is a tree, and $1$ is the root of the tree, 
              i.e. $\father 1 = 1$;
        \item for each object $A$ of $\mathcal{C}$ different from $1$, a morphism $p(A) \colon A \to \father A$
              called \textit{canonical projection of $A$}, we write it as $A \openrightarrow \father(A)$;
        \item for all objects $A, B$ in $\mathcal{C}$ with $A \lhd B$ and all morphisms $f \colon C \to A$, 
              an object $f^*B$ of $\mathcal{C}$ such that $C \lhd f^*B$, and a morphism $q(f, B) \colon f^*B \to B$
              in $\mathcal{C}$ such that
              \[
                \begin{tikzcd}
                    f^*B \arrow[r, "{q(f, B)}"] \arrow[d, open] & B \arrow[d, open] \\
                    C \arrow[r, "f"] & A
                \end{tikzcd}
              \]
              is a pullback in $\mathcal{C}$ that satisfies the following \textit{functoriality conditions}:
              \begin{enumerate}
                \item if $A$ and $B$ are objects in $\mathcal{C}$ such that $A \lhd B$, then 
                      \[
                        id_A^*B = B \quad \text{ and } \quad q(id_A, B) = id_B
                      \]
                \item whenever 
                      \[
                        \begin{tikzcd}
                            & & B \arrow[d, open] \\
                            A \arrow[r, "f"] & A' \arrow[r, "f'"] & A''
                        \end{tikzcd}
                      \]
                      in $\mathcal{C}$, then 
                      \[
                        (f' \circ f)^*B = f^*(f'^*B) 
                            \quad \text{ and } \quad 
                        q(f' \circ f, B) = q(f', B) \circ q(f, f'^*B)
                      \]
              \end{enumerate}
    \end{enumerate}
\end{definition}

We have that $R(U)$ is structured as a contextual category. The problem with $R(U)$ is that semantically equal 
but syntactically different rules are interpreted as different elements. Hence we will define an equivalence
relation between derived $T$- and $\in$- rules such that:
\[
        x_1 \colon A_1, \dots, x_n \colon A_n
        \Rightarrow
        A_{n+1} \type
    \equiv
        y_1 \colon B_1, \dots, y_m \colon B_m
        \Rightarrow
        B_{m+1} \type
\]
if, and only if, $m = n$ and for each $1 \leq i \leq n + 1$, \[
    \vdash x_1 \colon A_1, \dots, x_{i-1} \colon A_{i-1} \Rightarrow A_i = B_i[y_1 \mid\ x_1, \dots, y_{i-1} \mid\ x_{i-1}]
\]
and 
\[
        x_1 \colon A_1, \dots, x_n \colon A_n 
        \Rightarrow 
        t : A
    \equiv
        y_1 \colon B_1, \dots, y_m \colon B_m
        \Rightarrow
        s : B
\]
if, and only if, \[
        x_1 \colon A_1, \dots, x_n \colon A_n
        \Rightarrow
        A \type
    \equiv
        y_1 \colon B_1, \dots, y_m \colon B_m
        \Rightarrow
        B \type
\] and \[
    \vdash 
    x_1 \colon A_1, \dots, x_n \colon A_n 
    \Rightarrow t = s[y_1 \mid\ x_1, \dots, y_m \mid\ x_m] : A
\]

We define $C(U)$ as the homomorphic image of $R(U)$ induced by this equivalence relation.
This means that $C(U)$ is the category of equivalence classes of contexts and of equivalence
classes of realizations. Two contexts $x_1 \colon A_1, \dots, x_n \colon A_n$ and 
$y_1 \colon B_1, \dots, y_m \colon B_m$ are equivalent if, and only if, \[
        x_1 \colon A_1, \dots, x_{n-1} \colon A_{n-1}
        \Rightarrow
        A_{n} \type
    \equiv
        y_1 \colon B_1, \dots, y_{m-1} \colon B_{m-1}
        \Rightarrow
        B_{m} \type
\] and two realizations 
\begin{align*}
    \langle t_1, \dots, t_m \rangle 
        & \colon \langle x_1 \colon A_1, \dots, x_n \colon A_n \rangle
         \to  \langle y_1 \colon B_1, \dots, y_m \colon B_m \rangle \\
    \langle s_1, \dots, s_m \rangle
        & \colon \langle x'_1 \colon A'_1, \dots, x'_n \colon A'_n \rangle 
         \to \langle y'_1 \colon B'_1, \dots, y'_m \colon B'_m \rangle
\end{align*}
are equivalent if, and only if, for each $j$, $1 \leq j \leq n$,
\begin{multline*}
     x_1 \colon A_1, \dots, x_n \colon A_n
        \Rightarrow
        t_j \colon B_j[y_1 \mid\ t_1, \dots, y_{j-1} \mid\ t_{j-1}]
    \equiv \\ 
        x'_1 \colon A'_1, \dots, x'_n \colon A'_n
        \Rightarrow 
        s_j \colon B'_j[y'_1 \mid\ s_1, \dots, y_{j-1} \mid\ s_{j-1}]
\end{multline*}

It is easy to see that $C(U)$ is itself a contextual category.

\begin{definition}
    if $\mathcal{C}$ and $\mathcal{C}'$ are contextual categories, then a \textbf{contextual functor}
    $F \colon \mathcal{C} \to \mathcal{C}'$ is a functor $F \colon \mathcal{C} \to \mathcal{C}'$ such that:
    \begin{enumerate}
        \item $F(1) = 1$ and if $A \lhd B$ in $\mathcal{C}$, then $F(A) \lhd F(B)$ in $\mathcal{C}'$;
        \item for all object $A$ of $\mathcal{C}$, $F(p(A)) = p(F(A))$;
        \item for all $f$ and $B$ such that $f^*B$ is defined in $\mathcal{C}$,
            \[
                F(f^*B) = F(f)^*F(B) 
                    \quad \text{ and } \quad 
                F(q(f, B)) = q(F(f), F(B))
            \]
    \end{enumerate}
\end{definition}
$\mathbf{Con}$ is the category of contextual categories and contextual functors.

\section{Properties of Contextual Categories}

In this section we will state some properties about contextual categories that will be used later.

% generalized projection morphism
The first useful result is about the projections, although we defined projections on pairs $A \lhd B$, 
we might extend the notion to pairs $A \leq B$ defined as \[
    p(B, A) = p(X_1) \circ \dots, \circ p(X_n) \circ B
\] where $X_1, \dots, X_n$ is the unique sequence of objects in $\mathcal{C}$ such that
$A \lhd X_1 \lhd \dots \lhd X_n \lhd B$ in $\mathcal{C}$. Notation: $B \openrightarrow A$. % TODO: fix notation?

Let $\mathcal{C}$ be a contextual category and $A$ an arbitrary object of $\mathcal{C}$. Then we define 
the \textbf{relativization} of $\mathcal{C}$ to $A$, notes as $\mathcal{C}_A$, as the contextual category
where:
\begin{enumerate}
    \item the objects of $\mathcal{C}_A$ are the objects $B$ of $\mathcal{C}$ such that $A \leq B$.
    \item for objects $B, C$ in $\mathcal{C}_A$, a morphism from $B$ to $C$ over $A$ is a morphism, 
          i.e. a morphism from $B$ to $C$ in $\mathcal{A}$, is a morphism $f \colon B \to C$ in $\mathcal{C}$
          such that $p(C, A) \circ f = p(B, A)$. Composition of morphisms is inherited from $\mathcal{C}$.
    \item $\father_A(B) = \father(B)$ if $A < B$, and $\father_A(A) = A$.
    \item for objects $B > A$, $p_A(B) = p(B) \colon B \to \father(B)$.
    \item for $B, C, D$ objects of $\mathcal{C}_A$ with $\father_A(D) = C$ and morphisms $f \colon B \to C$ we put
          \[
                f^{*_A}D = f^*D \quad \text{ and } \quad q_A(f, D) = q(f, D)
          \]
\end{enumerate}

Now we will see that any morphism $f \colon B \to A$ in $\mathcal{C}$ induces a pullback functor 
$f^* \colon \mathcal{C}_A \to \mathcal{C_B}$ which preserves contextual structure.

\begin{lemma}
    let $\mathcal{C}$ be a contextual category, and $f \colon B \to A$ a morphism in $\mathcal{C}$,
    then for any object $C$ with $C \geq A$ we define an object $f^*C \geq B$ and a morphism 
    $q(f, C) \colon f^*C \to C$ by induction over $\level C - \level A$:
    \begin{enumerate}
        \item if $C = A$, \[
                f^*C = C \quad \text{ and } \quad q(f, C) = id_C
              \]
        \item if $C > A$, \[
                f^*C = q(f, \father(C))^*C
                    \quad \text{ and } \quad 
                q(f, C) = q(q(f, \father(C)), C)
              \]
    \end{enumerate}
    then if $C \geq A$, the pair $(p(f^*C, B), q(f, C))$ is a pullback cone over $(f, p(C, A))$.
\end{lemma}
\begin{theorem}
    Let $\mathcal{C}$ be a contextual category and $f \colon B \to A$ a morphism in $\mathcal{C}$, 
    by mapping an object $C$ of $\mathcal{C}_A$ to $f^*C$ in $\mathcal{C}_B$,
    and mapping $g \colon C \to D$ over $A$ to $f^*g \colon f^*C \to f^*D$ over $B$
    determined by
    \begin{align*}
        p(f^*D, B) \circ f^*g &= p(f^*C, B) \\
        q(f, D) \circ f^*g = q(f, C)
    \end{align*}
    we obtain the contextual functor $f^* \colon \mathcal{C}_A \to \mathcal{C}_B$
    called \textbf{reindexing along} $f$.
\end{theorem}
\begin{theorem} % TODO: see if it is used elsewhere or delete
    Let $\mathcal{C}$ be a contextual category and $f \colon B \to A$, $g \colon C \to B$ be morphisms in $\mathcal{C}$.
    Then, \[
        g^* \circ f^* = (f \circ g)^* 
            \quad \text{ and } \quad
        q(f \circ g, D) = q(f, D) \circ q(g, f^*D)
    \]
    for objects $D \geq A$.

\end{theorem}

\begin{definition}
    Let $\mathcal{C}$ be a contextual category, $A$ and $B$ be objects of $\mathcal{C}$ such that $A \lhd B$,
    the set of \textbf{sections} of $B$ is defined as 
    \begin{equation*}
        \Sections(B) \coloneqq \{ s \colon A \to B  \mid\ p(B) \circ s = \id_A\}
    \end{equation*}
\end{definition}
If $\level(B) = 1$, then 
\[
    \Sections(B) = \{ s \mid\ s \colon 1 \to B \}
\]
as any morphism $s \colon 1 \to B$ satisfies $p(B) \circ s = \id_1$ as 
$1$ is a terminal object. In this case, the notion of section
matches the notion of \textit{global element} of $B$.

\section{Relevance of the context category}

To see why this context category is relevant, we need to define interpretation fo a theory $U$
on a contextual category, and see that the interpretation into $C(U)$ is the 
initial object of the category of interpretations of $U$.

Given a contextual category $\mathcal{C}$ and a theory $U$, an \textbf{interpretation} of $U$ in $\mathcal{C}$
is a mapping $\mathcal{M}$ from the $T$- and $\in$-sequents of $U$ to the $\mathcal{C}$ such that
\begin{enumerate}
    \item if $R$ is a derived $T$-sequent \(
            \dfrac{
                x_1 \colon A_1, \dots, x_n \colon A_n
            }{
                A \text{ is a type}   
            }
        \) and we define sequents \\ 
        \(
            R_i \equiv \dfrac{
                x_1 \colon A_1, \dots, x_{i-1} \colon A_{i-1}
            }{
                A_i \text{ is a type}   
            }
        \) for $1 \leq i \leq n$, then $\mathcal{M}[R_1], \dots, \mathcal{M}[R_n], \mathcal{M}[R]$ are
        objects of $\mathcal{C}$ such that \[
            1 \lhd \mathcal{M}[R_1] \lhd \dots \lhd \mathcal{M}[R_n] \lhd \mathcal{M}[R]
        \]
        in $\mathcal{C}$.

    \item if $R_t$ is a derived $\in$-sequent \(
            \dfrac{
                x_1 \colon A_1, \dots, x_n \colon A_n   
            }{
                t \colon A
            }
        \) then $\mathcal{M}[R_t]$ is section of $\mathcal{M}[R]$.

    \item if $\Gamma \Rightarrow A = A'$ is a derived $T=$sequent, then if \[
            R \equiv \Gamma \Rightarrow A \type
            \quad \text{ and } \quad
            R' \equiv \Gamma \Rightarrow A' \type
        \] then $\mathcal{M}[R] = \mathcal{M}[R']$.
    \item if $\Gamma \Rightarrow t = t' \colon A$ is a derived $\in=$sequent, then if \[
            R_t \equiv \Gamma \Rightarrow t \colon A, \quad
            R_{t'} \equiv \Gamma \Rightarrow t' \colon A, \quad
            R \equiv \Gamma \Rightarrow A \type
        \] then, both $\mathcal{M}[R_t]$ and $\mathcal{M}[R_{t'}]$ are sections of $\mathcal{R}$
        and $\mathcal{M}[R_t] = \mathcal{M}[R_{t'}]$
\end{enumerate}

There is a canonical interpretation $\mathcal{M}$ of $U$ in $\mathcal{C}(U)$, defined as follows:
\begin{enumerate}
    \item if $R \equiv x_1 \colon A_1, \dots, x_n \colon A_n \Rightarrow A \type$ is a derived $T$-judgement, 
            then we define \[
                \mathcal{M}[R] = [x_1 \colon A_1, \dots, x_n \colon A_n, x_{n+1} \colon A]
            \]
    \item if $R_t \equiv x_1 \colon A_1, \dots, x_n \colon A_n \Rightarrow t \colon A$ is a derived $\in$-judgement,
            then we define \[
                \mathcal{M}[R_t] = [\langle x_1, \dots, x_n, t \rangle]
            \] a section of $\mathcal{M}[R]$.
\end{enumerate}

Now we shall define the concept of homomorphism between interpretations of a theory $U$.
Let $\mathcal{C}$, $\mathcal{C}'$ be two contextual categories, and let $\mathcal{M_C}$ and $\mathcal{M_{C'}}$
be two $U$-interpretations in $\mathcal{C}$ and $\mathcal{C}'$ respectively. Then, an interpretation 
homomorphism is a contextual functor $F \colon \mathcal{C} \to \mathcal{C}'$ such that:
\begin{enumerate}
    \item if $R$ is a $T$-rule, then $F \mathcal{M_C}[R] = \mathcal{M_{C'}}[R]$
    \item if $R_t$ is a $\in$-rule, then $F \mathcal{M_C}[R_t] = \mathcal{M_{C'}}[R_t]$
\end{enumerate}

So we can define the category $\mathbf{Con}_U$ as the category with pairs of contextual categories and $U$-interpretations
and their Homomorphisms. It is easy to see that the pair $(C(U), \mathcal{M})$ defined earlier is an initial object
of this category. This means that all the information of $U$ captured in any interpretation, is also captured by
this initial interpretation. 